\section{Scope, provenance, and statement of contributions}
\label{sec:scope}

This report continues the harmonic-tail reconstruction branch of the
\Repo\ polylogarithm manuscript.  The base repository was inspected at
commit~\cite{Repo}
\[
 \texttt{04bf92d507588a26abadbaee1beb1e7350c8bffb}
 \qquad (11\text{ October }2026).
\]
Three earlier continuations are the direct mathematical inputs.
The Bernoulli--Tornheim--digamma report established the shifted-Hankel
factorization and the reduction of each fixed smaller order $p=2a$ to
an $a\times a$ determinant.  The quartic spectral reconstruction
report proved effective eventual nonsingularity for every fixed $a$ and
computed the sufficient threshold $B_5=1920$.  The alternating-tail
report developed a five-square order-ten contraction and proved
nonsingularity outside one $2$-adic progression, but explicitly left
that progression and unrestricted order ten open.  The quartic report
then listed completion of $6\le b<1920$ as its first concrete research
target~\cite{BTD,QSR,ATSD}.

The present report resolves exactly that target.  It also records two
analytic consequences which are useful independently of the determinant
problem.  The first is a multiple-polylogarithm generating function for
two harmonic tails and its weighted-moment operator calculus.  The
second is a regularized depth-two multiple-polylogarithm generator for
integer-shift Stieltjes correlations.  These formulas were not found in
the inspected manuscript or continuation packages.  No claim of
priority over all external literature is made; the novelty claim in this
report is deliberately local: the formulas and proofs are new relative
to the audited \Repo\ corpus, and the exact order-ten completion answers
an open problem stated there.

\subsection{Main results}

The principal statements are the following.

\begin{enumerate}[label=\textbf{M\arabic*.}]
\item \textbf{Complete order-ten reconstruction.}
For every integer $b\ge6$, the rational coefficient matrix $M_{5,b}$
is nonsingular.  Hence the first $q=2b$ centered moments attached to
$(p,q)=(10,2b)$ have full formal rank $q$.

\item \textbf{One-prime finite certificate.}
For the finite range $6\le b\le1919$, reduction modulo
\[
 P=1{,}000{,}003
\]
certifies every determinant at once.  The $1914$ residues are all
nonzero.  Their canonical comma-separated encoding has digest
\[
 \texttt{3977e54f5b43db43e8f080c32f2e5b22beac54412cca32ff5d03d84975111836}.
\]
The certificate is independently regenerated from the Bernoulli
recurrence and exact modular Gaussian elimination.

\item \textbf{Ordinary reconstruction identities.}
For each $b\ge6$ and $0\le k<b$, explicit rational polynomials
$P_{b,k}$ and $Q_{b,k}$ satisfy
\[
 \sum_{n=0}^{\infty}
 \left(P_{b,k}(n+\tfrac12)
       [\zeta(10)-H_n^{(10)}][\zeta(2b)-H_n^{(2b)}]
       -Q_{b,k}(n+\tfrac12)\right)=\zeta(2k),
\]
where $\zeta(0)=-1/2$.  The same identity is a product-polygamma
formula after replacing each tail by a polygamma derivative.

\item \textbf{Two-tail polylogarithmic generator.}
For $p,q>1$ and $|z|<1$, the ordinary generating function of
$\zeta(p,n+1)\zeta(q,n+1)$ is a divided difference of depth-two
multiple polylogarithms.  Polynomially weighted centered moments are
obtained by the single Euler operator $z\,d/dz+1/2$.

\item \textbf{Stieltjes correlation generator.}
A bivariate pole-subtracted depth-two polylogarithm kernel is holomorphic
at $(\alpha,\beta)=(0,0)$, and its Taylor coefficients generate all
products $\gamma_r(n+1)\gamma_s(n+1)$.  The theorem includes explicit
anchored primitives and special values.
\end{enumerate}

\subsection{Proof-status conventions}

A statement marked \emph{inherited} is quoted from an existing \Repo\
report and used as an input.  A statement marked \emph{new} is proved
here.  A numerical check is never used as a proof gate.  The finite
certificate operates in an exact prime field; the high-precision check
of the explicit $(10,16)$ identity is only a diagnostic.  Formal rank
means rank in the rational coordinate space specified by the source
construction.  It does not assert arithmetic independence of the
numbers $\zeta(2k)$ or of any periods obtained after numerical
evaluation.

\subsection{Why a finite certificate is a proof here}

The eventual theorem converts the order-ten conjecture into a finite
problem before any computation begins.  The retained determinant has
fixed size five, every entry is an explicitly displayed rational
function of $b$ and finitely many Bernoulli numbers, and all denominators
are controlled.  Reduction modulo a prime which divides none of those
denominators is an exact ring homomorphism.  A nonzero modular
determinant therefore proves that the original rational determinant is
nonzero.  The certificate is not evidence extrapolated from samples;
it is an exhaustive exact verification of the complete finite set left
by the analytic theorem.
